\vspace{-1em}
\section{Discussion and Conclusion}
\vspace{-1em}

We introduce OpenShape, a novel approach for learning scalable and generalizable multi-modal joint representations for 3D shapes. OpenShape representations effectively capture a wide range of semantic and visual concepts, enabling superior capabilities for open-world 3D shape recognition. By aligning OpenShape with CLIP's embedding space, our shape embeddings can be integrated with off-the-shelf CLIP-based models for various cross-modality applications. Moving forward, there are several directions worth further exploration: (a) More 3D data. While we utilized 876k 3D shapes during training, this is still quite limited compared to the 2D counterparts. We hope that our work inspires future investments in more resources to build even more powerful 3D representations. (b) Part-level information. Our current shape representations mainly focus on global semantic and visual features, and it would be beneficial to add more part-level supervision during training. (c) Sim-to-real domain gap. Our model is mainly trained on synthetic data, and it's challenging but crucial to explore explicit designs for reducing the domain gap with real-world shapes.

